\scanpage{61}
\mspage{319}

Let us now consider the homomorphism \unsure{(3)}
\[
\tag{18}
\left\{
\begin{aligned}
\widehat\varphi:\quad &\widehat\Pi_{0,3}\longrightarrow
  \widehat{\mathfrak S}_{0,3}^{+}\\
&\ell_0\longmapsto\ell'_0=\sigma_\infty\rho=\varepsilon_0\\
&\ell_1\longrightarrow\ell'_1=\sigma_\infty\\
&\ell_\infty\longmapsto\ell'_\infty=\rho^{-1}
\end{aligned}
\right.
\]
and let us write that it commutes, \unsure{as} an outer homomorphism, with $u$.
This therefore means that there exists a
\[
\tag{19}
\lambda\in\widehat{\mathfrak S}_{0,3}^{+}
\]
such that we have
\[
\tag{20}
\widehat\varphi(\unsure{ux})=\lambda\,u\widehat\varphi(x)\,\lambda^{-1}
\qquad \forall x\in\widehat\Pi_{0,3},
\]
and it suffices to express this condition for $x$ equal to one of the two
generators $\ell_0,\ell_1$ of $\widehat\Pi_{0,3}$

For $x=\ell_0$ relation (20) is written
\[
\tag{21}
u(\ell'_0)=\operatorname{int}(\lambda^{-1}\widehat\varphi(g_0))(\ell_0'{}^{p})
\]
which, compared with (17), gives
\[
\operatorname{int}(g_0^{-1}\lambda^{-1}\widehat\varphi(g_0))(\ell_0'{}^{p})
=\ell_0'{}^{p}\quad\text{hence}
\]
\[
\deleted{(22)}\qquad
g_0^{-1}\lambda^{-1}\widehat\varphi(g_0)=\ell_0'{}^{\beta}
\qquad(\beta\in\widehat{\mathbb Z})
\]
or again
\[
\tag{22}
\boxed{\lambda=\widehat\varphi(g_0)\ell_0'{}^{-\beta}g_0^{-1}}
\inserted{=\widehat\varphi(g_0h_0^{-1})=\widehat\varphi(\sigma_\infty(g_0))}
\]
\margin{$\searrow$ \unsure{cf.} below (24 bis)}

For $x=\ell_1$ relation (20) is written

\scanpage{62}

\[
\tag{23}
u(\ell'_1)=\operatorname{int}(\lambda^{-1}\widehat\varphi(g_1))(\ell_1'{}^{p})
\]
\margin{where $(g_1=\sigma_\infty(g_0)=\sigma_\infty g_0\sigma_\infty\unsure{)}$}
or again
\[
\sigma_\infty=\operatorname{int}(\lambda^{-1}\widehat\varphi(g_1))(\sigma_\infty)
\]
\margin{For $\sigma_\infty^{p}=\sigma_\infty$ since
\inserted{$\sigma_\infty^2=1$ and} \unsure{the image} of $p$ in
$(\mathbb Z/2\mathbb Z)^*$ is $1$!}
or again $\sigma_\infty$ commutes with
\[
\lambda^{-1}\widehat\varphi(g_1)
\overset{\inserted{\text{cf. (22)}}}{=}
g_0\ell_0'{}^{\beta}\widehat\varphi
 \bigl(\underbrace{g_0^{-1}\sigma_\infty(g_0)}_{h_0^{-1}}\bigr)
=g_0\ell_0'{}^{\beta}\widehat\varphi(h_0^{-1})
\]
or again with its inverse $\widehat\varphi(h_0)\ell_0'{}^{-\beta}g_0^{-1}$
which is also written
\[
\underbrace{\sigma_\infty
 \bigl(\widehat\varphi(h_0).\ell_0'{}^{-\beta}.g_0^{-1}\bigr)}
=\widehat\varphi(h_0).\ell_0'{}^{-\beta}.g_0^{-1}
\]
\[
\sigma_\infty(\widehat\varphi(h_0)).
\underbrace{\varepsilon_1^{-\beta}}_{(\rho\sigma_\infty)^{-\beta}}.
\sigma_\infty(g_0)^{-1}
=\widehat\varphi(h_0)\ell_0'{}^{-\beta}g_0^{-1}
\]
\[
\sigma_\infty(\widehat\varphi(h_0))
 (\underbrace{\rho\sigma_\infty}_{\varepsilon_1})^{-\beta}h_0
=\widehat\varphi(h_0).\varepsilon_0^{-\beta}
\]
\[
\tag{24}
\boxed{
\overbrace{\varepsilon_1^{-\beta}h_0\varepsilon_0^{\beta}}
 ^{\inserted{\sigma_\infty(g_0\varepsilon_0^{\beta})^{-1}(g_0\varepsilon_0^{\beta})}}
=\sigma_\infty(\widehat\varphi(h_0))^{-1}\widehat\varphi(h_0)
}
\]
\margin{\unsure{Correct?} $\searrow$}
\margin{\uline{N.B.} If $g_0$ is replaced by
\inserted{$g'_0=$} $g_0\ell_0^{\lambda}$ $(\lambda\in\widehat{\mathbb Z})$,
$h_0$ is replaced by $h'_0=\ell_1^{-\lambda}\unsure{h_0|g_0}\ell_0^{\lambda}$,
and we \unsure{find} that $\beta$ \unsure{is} \unsure{''} replaced by
$\beta'=\beta-3\lambda$.
This shows \inserted{?} that we can in a \uline{unique way}
\illegible{1} \uline{normalize} the choice of $g_0$, in such a way that we have
$\beta\in\{0,1,2\}$.
(It would be necessary to elucidate whether $\beta$ actually takes all
\unsure{three} values modulo $3$\ldots)}
\margin{\uuline{N.B.} This equation is equivalent to the
\uline{simpler} equation
\[
\text{(24 bis)}\qquad
\boxed{\widehat\varphi(h_0)=\sigma_\infty^{\unsure{\varepsilon}}g_0\varepsilon_0^{\beta}}
\]}
\[
\unsure{(25)}\qquad
\left\{
\begin{aligned}
&h_0\text{ given by (8) in terms of }g_0\\
&\varepsilon_0=\ell'_0=\sigma_\infty\rho
   &&(\varepsilon_0^2=\ell_0)\\
&\varepsilon_1=\rho\sigma_\infty=\sigma_0\rho
   &&(\varepsilon_1^2=\ell_1)
\end{aligned}
\right.
\]

Setting, for $\deleted{\unsure{g}}\ \mu\in
\widehat{\mathfrak S}_{0,3}^{+}$,
\[
\tag{26}
\widetilde\mu=\sigma_\infty(\mu)^{-1}\mu
=\sigma_\infty\mu^{-1}\sigma_\infty\mu=[\sigma_\infty,\mu]
\]
(so that $h_0=\widetilde g_0$) the formula \unsure{(24)|(25)} takes the form
\[
\tag{27}
\varepsilon_1^{-\beta}\unsure{\widetilde g_0}\varepsilon_0^\beta
=\bigl(\widehat\varphi(\widetilde g_0)\bigr)^{\sim}.
\]

\scanpage{63}
\mspage{320}

In summary:

\uline{Proposition} Let us consider an outer automorphism $u$ of the group with loops
$\widehat\Pi_{0,3}$, realized by $g_0\in\widehat\Pi_{0,3}$
\inserted{modulo $L_0$,}
and $p\in\widehat{\mathbb Z}^{*}$,
(and $\gamma\in\widehat{\mathbb Z}^{*}$ uniquely determined
\unsure{by terms|in terms} of $g_0$ and $p$)
by formulas (11), by means of an actual automorphism
\inserted{(still denoted $u$)} which commutes with $\sigma_\infty$ ---
the data $g_0,p$ (and $\gamma$) \unsure{being} \inserted{therefore}
required to satisfy conditions (9) and (10) above,
plus the generation condition (11 bis).
\inserted{(We extend $u$ \unsure{to} an automorphism of
$\widehat{\mathfrak S}_{0,3}^{+}$
by $\operatorname{int}(u)$, cf. (15))}
For the outer homomorphism
$\widehat\varphi:\widehat\Pi_{0,3}\longrightarrow
\widehat{\mathfrak S}_{0,3}^{+}$
defined by (18) to commute with the outer action of $u$,
it is necessary and sufficient that we can find a $\beta\in\widehat{\mathbb Z}$,
such that we have (24), or what amounts to the same, (27),
\margin{\unsure{better:} \uuline{24 bis}}

This condition \inserted{(for $\beta$ fixed)} is also equivalent to formula
(20) for $x=\ell_\infty$ ($\lambda$ \unsure{being} always given by (22)),
\deleted{\unsure{(2}} i.e. to the formula
\[
\tag{28}
u(\ell'_\infty)=\operatorname{int}(\lambda^{-1}\widehat\varphi(g_\infty))
 (\ell_\infty'{}^p),\qquad\text{where}
\]
\[
\tag{29}
g_\infty=g_0\ell_0^{-\gamma}\rho^{-1}(h_0)
\qquad(\text{and }\lambda\text{ given by (22)}).
\]

\scanpage{64}

formula which thus makes precise the \unsure{presumptive} formula (18). It is also written, more explicitly
\[
\rho^{-1}g^{-1}
=\operatorname{int}\bigl(g_0\varepsilon_0^{\beta}
\widehat\varphi(\ell_0^{\gamma}\rho^{-1}(h_0))\bigr)
(\rho^{-p})
\]
\inserted{\[
g_0\varepsilon_0^{\beta+\gamma}\widehat\varphi(\rho^{-1}(h_0))
\]}
i.e. (taking inverses, and making $g$ explicit)
\[
(30)\qquad
\boxed{
\deleted{\illegible{1}}\inserted{h_0^{-1}\ell_1^{\gamma}\rho}
=\operatorname{int}\bigl(
\deleted{\unsure{g_0}}\varepsilon_0^{\deleted{\unsure{\beta+\gamma}}\inserted{\alpha}}
\widehat\varphi(\rho^{-1}(h_0))\bigr)
(\rho^{p\deleted{\unsure{-p}}})
}
\]
\margin{\uline{N.B.} We put $\alpha=\beta+\gamma$. The existence \unsure{of} $\beta\in\widehat{\mathbb Z}$ satisfying (24), is equivalent to that \unsure{of} $\alpha\in\widehat{\mathbb Z}$ satisfying (30).}

The equivalence of (24) and (30) may seem a little strange, owing to the fact \inserted{for example} that in (24) the exponent $p$ does not enter explicitly, but only $h_0$, whereas formula (30) brings in the value of $p\in\widehat{\mathbb Z}^{*}$ modulo $3$ (which gives $\rho^p=\rho$ or $\rho^p=\rho^{-1}$). But in any event, whether it be in the form (24) \unsure{where|or} in the form (30), it does seem that the condition in question is highly non-tautological!

To do things properly, it would be necessary to \unsure{first} make explicit the congruential conditions of \S~39 in terms of $g_0,p$.

\scanpage{65}
\mspage{321}

I would like to \unsure{compare} the outer automorphisms of $\widehat\Pi_{0,3}$, and those of $\widehat{\mathfrak S}_{0,3}^{+}$, taking account of the exact sequence
\[
(31)\qquad
\diagram{\begin{tikzcd}[ampersand replacement=\&]
1\arrow[r,no head]\&\widehat\Pi_{0,3}\arrow[r,no head]\&
\widehat{\mathfrak S}_{0,3}^{+}\arrow[r]\arrow[d,draw=none,"{\inserted{\unsure{\wr}}}" description]\&
\mathfrak S_{3}\arrow[r,no head]\&1\\
\&\&\inserted{\widehat{\Pi'}_{0,3}}
\end{tikzcd}}
\]

\deleted{\uline{Proposition}} \inserted{Let us note} that every outer automorphism \inserted{$u$} of $\widehat{\mathfrak S}_{0,3}^{+}$ which \unsure{leaves invariant} the normal subgroup $\widehat\Pi_{0,3}$, and induces on the quotient $\mathfrak S_{3}$ a trivial outer automorphism, can be written as obtained from an automorphism \inserted{$u_0$} of $\widehat{\mathfrak S}_{0,3}^{+}$, \unsure{stabilizing} $\widehat\Pi_{0,3}$, and inducing the \uline{identity} on the quotient $\mathfrak S_{3}$ --- and $u_0$ is then determined modulo composition with an element $\operatorname{int}(g)$, where $g\in\widehat{\mathfrak S}_{0,3}^{+}$ \unsure{is} such that its image in $\mathfrak S_{3}$ is central, i.e. $=1$, i.e. $g\in\widehat\Pi_{0,3}$. Thus
\[
(32)\qquad
\begin{gathered}
\text{group of outer automorphisms}\\
\text{of $\widehat{\mathfrak S}_{0,3}^{+}$ which}\\
\text{induce on $\mathfrak S_{3}$ a}\\
\text{trivial outer automorphism (hence \unsure{leave invariant}}\\
\widehat\Pi_{0,3})
\end{gathered}
\simeq
\begin{gathered}
\text{group of}\\
\text{automorphisms of}\\
\widehat{\mathfrak S}_{0,3}^{+}\text{ which}\\
\text{induce}\\
\text{the identity on}\\
\text{its quotient $\mathfrak S_{3}$,}\\
\text{\uline{modulo} the sub-}\\
\text{group of}\\
\text{inner automorphisms}\\
\text{by $g\in\widehat\Pi_{0,3}$}
\end{gathered}
\]

\scanpage{66}

Is \inserted{such} an automorphism $u_0$ determined by its restriction to $\widehat\Pi_{0,3}$? In other words, an automorphism of $\widehat{\mathfrak S}_{0,3}^{+}$ \deleted{which} ($\simeq\mathrm{Sl}(\unsure{2|1},\mathbb Z)^{\unsure{\prime}\wedge}$) which induces the identity on $\widehat\Pi_{0,3}$ and on $\mathfrak S_{3}$, is it trivial? It suffices to prove that its restriction to the inverse image \inserted{$\widehat\Pi_{\unsure{0,\infty|0,3}}\cdot\{1,\unsure{\sigma_\infty}\}$} of $\{1,\dot\sigma_\infty\}\subset\mathfrak S_{3}$ is \unsure{trivial|trivial} (by symmetry this will also be the case on replacing $\dot\sigma_\infty$ by $\dot\sigma_0,\dot\sigma_1$, now these three subgroups generate $\widehat{\mathfrak S}_{0,3}^{+}$. We will have $f(\sigma_\infty)=\sigma_\infty g$, with $g\in\widehat\Pi_{0,3}$ [and we must have $\sigma_\infty g\sigma_\infty g=1$ i.e.
\[
g\cdot\sigma_\infty(g)=1
\]
]
moreover \deleted{for $x\in\Pi_{0,3}$ \unsure{we have}} \inserted{we will have to have that}
\[
\operatorname{int}(\sigma_\infty g)\vert\widehat\Pi_{0,3}
=\deleted{\unsure{f(x\sigma_\infty)}}\operatorname{int}(\sigma_\infty)\vert\widehat\Pi_{0,3}
\]
whence \deleted{$g\in$}
\[
\operatorname{int}(g)\vert\widehat\Pi_{0,3}=\mathrm{id}
\quad\text{i.e.}\quad
g\in\unsure{\operatorname{Centre}}(\widehat\Pi_{0,3})
\quad\text{i.e.}\quad g=1,
\]
OK.

But under what condition does an automorphism of $\widehat\Pi_{0,3}$ \deleted{can} extend to an automorphism of $\widehat{\mathfrak S}_{0,3}^{+}$ which \deleted{is}\inserted{\unsure{induces}} the identity on $\mathfrak S_{3}$? We see at once that it \uline{is necessary} for this that the outer automorphism of $\widehat\Pi_{0,3}$ which it induces commute with the action of $\mathfrak S_{3}$.

\scanpage{67}
\mspage{\unsure{322|332}}

This is sufficient, for let $\Gamma\inserted{{}\subset\mathrm{Aut\,ext}(\widehat{\Pi}_{0,3})}$ be the group of outer
automorphisms of $\widehat{\Pi}_{0,3}$ which have this property (we are not
concerned for the moment with the
loop structure of $\widehat{\Pi}_{0,3}$), we have
$\Gamma\cap\mathfrak S_3=(1)$ (identifying $\mathfrak S_3$ with a
\unsure{subgroup} of $\mathrm{Aut\,ext}\,\widehat{\Pi}_{0,3}$) for \deleted{\unsure{$\mathfrak S$}} the center
of $\mathfrak S_3$ is trivial, hence
\margin{this is \unsure{evident|evidently}}
\[
\deleted{\unsure{\Gamma\times}}\ \mathfrak S_3\times\Gamma\hookrightarrow\mathrm{Autext}(\widehat{\Pi}_{0,3}),
\]
and we thus have an extension structure
\[
1\unsure{\longrightarrow}\widehat{\Pi}_{0,3}\longrightarrow\mathcal E\longrightarrow\mathfrak S_3\times\Gamma\longrightarrow1
\tag{33}
\]
(for the center of $\widehat{\Pi}_{0,3}$ is trivial), or again
an extension structure
\[
1\longrightarrow\widehat{\mathfrak S}_{0,3}^{+}\unsure{\longrightarrow}\mathcal E\longrightarrow\Gamma\longrightarrow1
\tag{34}
\]
which defines an outer action of $\Gamma$
on $\widehat{\mathfrak S}_{0,3}^{+}$ --- which is the outer action
sought. Thus

\uline{Proposition} the group of outer automorphisms of $\widehat{\Pi}_{0,3}$ which
\unsure{commute} with $\mathfrak S_3$ is canonically isomorphic to the group
of outer automorphisms of $\widehat{\mathfrak S}_{0,3}^{+}$ which induce
on $\mathfrak S_3$ the trivial outer automorphism.
\margin{But \uline{every} automorphism of $\mathfrak S_3$ is inner!}

\scanpage{68}

\uline{Corollary} the group of automorphisms of $\widehat{\Pi}_{0,3}$ which
commute as outer actions with $\mathfrak S_3$, is isomorphic
to the group of automorphisms of $\widehat{\mathfrak S}_{0,3}^{+}$ which
induce on $\mathfrak S_3$ the trivial automorphism (by
restriction of the latter)

\uuline{Corollary}
\[
\diagram{\begin{tikzcd}[ampersand replacement=\&,column sep=small,row sep=large]
1 \arrow[r,no head]
\& \widehat{\Pi}_{0,3} \arrow[r]
\& \mathrm{Aut}\left(\widehat{\Pi}_{0,3}\,\middle|\,
  \begin{array}{l}\mathrm{comm.}\\\mathrm{ext.}^{\unsure{t}}\\\text{with }\mathfrak S_3\end{array}\right) \arrow[r]
\& \mathrm{Autext}\left(\widehat{\Pi}_{0,3}\,\middle|\,
  \begin{array}{l}\mathrm{comm.}\\\text{with }\mathfrak S_3\end{array}\right) \arrow[r]
\& 1
\\
1 \arrow[r,no head]
\& \widehat{\Pi}_{0,3} \arrow[r,no head] \arrow[u,phantom,"{\text{''}}" description]
\& \underbrace{\mathrm{Aut}\left(\widehat{\mathfrak S}_{0,3}^{+}\,\middle|\,
  \begin{array}{l}\text{identity on the}\\\text{quotient }\mathfrak S_3\end{array}\right)}_{\mathcal E^{!}}
  \arrow[r] \arrow[u,"{\sim\ \text{\unsure{restriction} to }\widehat{\Pi}_{0,3}}"]
\& \mathrm{Autext}\left(\widehat{\mathfrak S}_{0,3}^{+}\,\middle|\,
  \begin{array}{l}\text{\deleted{\unsure{finite}}\inserted{\unsure{trivial}}}\\\text{on the}\\\text{quotient }\mathfrak S_3\end{array}\right)
  \arrow[r] \arrow[u,"\sim"]
\& \unsure{1}
\end{tikzcd}}
\tag{35}
\]

\deleted{\uline{Corollary}
Let $u\in\mathrm{Aut}\deleted{\mathrm{ext}}(\widehat{\mathfrak S}_{0,3}^{+})$, \unsure{inducing}
\unsure{Let} \illegible{1} $u\in\mathrm{Aut}(\widehat{\mathfrak S}_{0,3}^{+})$ and $\unsure{u_0}\unsure{=|:=}u|\widehat{\Pi}_{0,3}$
\margin{in $\mathcal E$}
the trivial \inserted{outer} automorphism on $\mathfrak S_3$. We thus have
\unsure{commuting} \illegible{2}
\[
u(\sigma_\infty)=g\sigma_\infty\quad\text{with}\quad g\in\widehat{\Pi}_{0,3}\quad\text{\unsure{well} determined},
\]
we have \illegible{1} that $\sigma_\infty$
and $(g\sigma_\infty)^2=1$ i.e. $g\sigma_\infty g\sigma_\infty=1$ i.e. $g\sigma_\infty g$.}

This extension $\mathcal E^{!}$ of $\Gamma$ by $\widehat{\Pi}_{0,3}$ is
split\inserted{\unsure{noncanonically}}, the \unsure{splitting|splitting} subgroup consists
of the automorphisms of $\widehat{\Pi}_{0,3}$ which commute with $\sigma_\infty$,
or again the automorphisms of $\widehat{\mathfrak S}_{0,3}^{+}$ which
\deleted{\unsure{commute}}\inserted{fix} \unsure{with} $\sigma_\infty$ (and \inserted{which fix} $\rho\ \mathrm{mod}\ \widehat{\Pi}_{0,3}$)
\margin{comes from the fact that there \deleted{\illegible{2}} \inserted{\unsure{exists}} \deleted{\illegible{1}} \unsure{a} \illegible{2} \inserted{$\widehat{\Pi}_{0,3}$-\unsure{conjugation}} \illegible{1} of order $2$ above $\sigma_\infty$ \ldots}
\[
\left\{
\begin{aligned}
u(\sigma_\infty)&=\sigma_\infty\\
u(\rho)&=g.\rho
\end{aligned}
\right.
\qquad(g\in\widehat{\Pi}_{0,3})
\tag{36}
\]

\scanpage{69}
\mspage{323}

\unsure{writing} that $u$ is an \unsure{endomorphism} (for we \unsure{recall}
that $\widehat{\mathfrak S}_{0,3}^{+}$ is defined by the generators $\sigma_\infty,\rho$
with the relations
\begin{equation}
\sigma_\infty^2=\rho^3=1
\tag{37}
\end{equation}
\unsure{we find} the condition
\begin{equation}
(g\rho)^3=1\qquad\text{i.e.}\qquad g\,\rho(g)\,\rho^2(g)=1.
\tag{38}
\end{equation}

\margin{moreover it \unsure{is necessary} in order that $\sigma_\infty$ and $g\rho$
\unsure{generate} $\unsure{\widehat\Pi_{0,3}}$, with only relations
$\sigma_\infty^{\unsure{2}}=(g\rho)^3=1$.}

\uline{Corollary}. Let $u\in\mathcal E^{!}$ be an automorphism of
$\widehat{\mathfrak S}_{0,3}^{+}$ which induces the identity on $\mathfrak S_3$.
Equivalent conditions:

a) $\exists\ p\in\widehat{\mathbb Z}^{*}$ such that
$u(\underset{\varepsilon_0}{\uline{\ell'_0}})$ conjugate in
$\widehat{\mathfrak S}_{0,3}^{+}$ to $\ell_0'{}^p$
(this $p$ is then unique\inserted{, and we can conjugate by an \unsure{element} of $\widehat\Pi_{0,3}$}).

b) $\exists\ p\in\widehat{\mathbb Z}^{*}$ such that $u(\ell_0)$ is conjugate in
$\widehat\Pi_{0,3}$ to $\ell_0^p$ (this $p$ is then unique).

c) $u_0\overset{\mathrm{dfn}}{=}u\mid\widehat\Pi_{0,3}$ respects the loop structure of $\widehat\Pi_{0,3}$.

Moreover, the $p$ in a) and b) is the same.

\uline{\unsure{Proof}}. If $u(\ell'_0)=\operatorname{int}(g_0)(\ell_0'{}^p)$, with
$g_0\in\widehat{\mathfrak S}_{0,3}^{+}$,
squaring we \unsure{will have} $u(\ell_0)=\operatorname{int}(g_0)(\ell_0^p)$,
this shows that the image $\dot g_0$ of $g_0$ in $\mathfrak S_3$
\unsure{leaves invariant} $0$ i.e. is \deleted{\illegible{1}} $1$ or $\dot\sigma_0$,
\deleted{therefore} if it is $1$ we do have b), otherwise we \unsure{put|\deleted{know q}}
$g_0=h\sigma_0$ \inserted{($h\in\widehat\Pi_{0,3}$)}
\deleted{\unsure{whence}}
\[
u(\ell_0)=\operatorname{int}(h)\,
\underset{(\sigma_0(\ell_0))^p}{\uline{\sigma_0(\ell_0^p)}}
\]
now we know that $\sigma_0(\ell_0)=\operatorname{int}(a)\ell_0$
\inserted{(with $a\in\widehat\Pi_{0,3}$)},
\inserted{whence $\sigma_0(\ell_0)^p=\operatorname{int}(a)\ell_0^p$},
thus $u(\ell_0)=\operatorname{int}(ha)(\ell_0^p)$, whence b) with the same $p$.

\scanpage{70}

Since we know the uniqueness of $p$ in b), in a) too. If we have
\begin{equation}
u_0(\ell_0)=\operatorname{int}(g_0)(\ell_0^p)
\qquad (g_0\in\widehat\Pi_{0,3})
\tag{39}
\end{equation}
we deduce from this, writing
\[
\deleted{\rho u_0\rho^{-1}(\ell_0)\unsure{=}}
\qquad
\rho u_0\rho^{-1}=\operatorname{int}(g^{-1})u_0
\qquad (\unsure{a|g}\in\widehat\Pi_{0,3})
\]
\[
\deleted{u_0 \ell_1}
\]
\[
u_0(\ell_1)=\operatorname{int}(\unsure{a|g})
\underbrace{\rho\underbrace{u_0\underbrace{\rho^{-1}(\ell_1)}_{\ell_0}}_{
\operatorname{int}(g_0)(\ell_0^p)}}_{
\operatorname{int}(\rho(g_0))(\ell_1^p)}
=\operatorname{int}(g\rho(g_0))(\ell_1^p)
\]
i.e.
\begin{equation}
u_0(\ell_1)=\operatorname{int}(g_1)(\ell_1^p)
\qquad\text{with}\quad g_1=g\rho(g_0)
\tag{40}
\end{equation}
and likewise
\begin{equation}
\unsure{u_0|u}(\ell_\infty)=\operatorname{int}(g_\infty)(\ell_\infty^p)
\qquad\text{with}\quad
\begin{aligned}
g_\infty&=g\rho(g_1)\\
&=g\rho(g)\rho^2(g_0)
\end{aligned}
\tag{41}
\end{equation}
thus b) $\Rightarrow$ c). Finally \unsure{let us prove} c) $\Rightarrow$ a).

We have by hypothesis
\[
\begin{array}{rcl}
u_0(\ell_0)&=&\operatorname{int}(g)\ell_0^p\\
\text{''}&&\text{''}\\
u(\varepsilon_0)^2&&
\operatorname{int}(g)(\unsure{\ell|\varepsilon}_0^{\unsure{2p|p}})^2
=(\operatorname{int}g(\unsure{\ell_0^p|\varepsilon_0^p}))^2
\end{array}
\]
I say that we then have
\[
u(\varepsilon_0)=\operatorname{int}(g)(\varepsilon_0^p).
\]
It suffices to \deleted{see}\inserted{\unsure{recall}} that
\deleted{\unsure{every}}\inserted{the only} \unsure{element} of
$\widehat{\mathfrak S}_{0,3}^{+}$ with square equal to $\ell_0$ is $\varepsilon_0$\ldots\quad OK

\margin{\unsure{check}}
